\addcontentsline{toc}{subsubsection}{Solution 61. The Hidden Entropy Function}

\begin{dlsummary}[title={\bfseries Solution 61. The Hidden Entropy Function},breakable]

\begin{proof}
The unique solution is
\[
\boxed{f(x)=x\ln x-x+1\qquad(x>0).}
\]

\paragraph{1. Find the forced normalisation.}

Putting \(x=y=1\),
\[
f(1)\ge2f(1).
\]
Nonnegativity therefore gives
\[
f(1)=0.
\]

\paragraph{2. Absorb the correction term.}

Define
\[
g(x)=\frac{f(x)+x-1}{x}.
\]
A direct rearrangement of the original inequality gives
\[
\boxed{g(xy)\ge g(x)+g(y).} \tag{1}
\]
Also, \(g(1)=0\), and nonnegativity of \(f\) implies
\[
\boxed{g(x)\ge1-\frac1x.} \tag{2}
\]

The problem is now to sharpen the elementary bound (2).

\paragraph{3. Subdivide multiplicatively.}

For every positive integer \(n\), repeated application of (1) gives
\[
g(x)\ge n\,g(x^{1/n}).
\]
Using (2),
\[
g(x)\ge n\left(1-x^{-1/n}\right).
\]
The elementary numerical limit
\[
\lim_{n\to\infty}n\left(1-x^{-1/n}\right)=\ln x
\]
therefore yields
\[
\boxed{g(x)\ge\ln x\qquad(x>0).} \tag{3}
\]

No continuity of \(g\) has been used: the limit concerns only the explicit expression on the right.

\paragraph{4. Use reciprocals to obtain the reverse bound.}

Apply (1) to \(x\) and \(1/x\):
\[
0=g(1)\ge g(x)+g(1/x).
\]
By (3),
\[
g(x)\le-g(1/x)\le-\ln(1/x)=\ln x.
\]
Thus \(g(x)=\ln x\). Undoing the normalisation,
\[
f(x)=xg(x)-x+1=x\ln x-x+1.
\]

\paragraph{5. Verify.}

First, the standard inequality
\[
\ln u\le u-1
\]
applied to \(u=1/x\) gives
\[
x\ln x-x+1\ge0.
\]
So the required range is satisfied.

Next, direct expansion gives
\[
\begin{aligned}
&f(xy)-xf(y)-yf(x)\\
&\qquad=xy-x-y+1\\
&\qquad=(x-1)(y-1).
\end{aligned}
\]
Hence the original inequality holds with equality.
\end{proof}
\end{dlsummary}

\begin{takeaways}
\begin{itemize}
\item \textbf{Choose the normalisation strategically:} Adding \(x-1\) and dividing by \(x\) removes both the weights and the correction term.
\item \textbf{Use the range as structural information:} Nonnegativity becomes the lower bound that drives the entire argument.
\item \textbf{Repeated roots sharpen a crude estimate:} Multiplicative subdivision turns \(1-1/x\) into \(\ln x\).
\item \textbf{Reciprocals force exactness:} Once the global lower bound is established, the upper bound follows immediately.
\end{itemize}
\end{takeaways}
